\documentclass[11pt,a4paper]{article}

\usepackage[utf8]{inputenc}
\usepackage[T1]{fontenc}
\usepackage{lmodern}
\usepackage[margin=2.5cm]{geometry}
\usepackage{amsmath,amssymb,amsthm}
\usepackage{graphicx}
\usepackage{booktabs}
\usepackage{xcolor}
\usepackage{hyperref}
\usepackage{listings}
\usepackage{array}
\usepackage{multirow}
\usepackage{caption}
\usepackage{subcaption}
\usepackage{enumitem}

\hypersetup{
    colorlinks=true,
    linkcolor=blue!60!black,
    citecolor=blue!60!black,
    urlcolor=blue!60!black
}

\lstset{
    basicstyle=\ttfamily\small,
    breaklines=true,
    frame=single,
    backgroundcolor=\color{gray!8},
    keywordstyle=\color{blue!70!black},
    commentstyle=\color{green!50!black},
    stringstyle=\color{red!60!black}
}

% Notation shortcuts
\newcommand{\bsigma}{\boldsymbol{\sigma}}
\newcommand{\beps}{\boldsymbol{\varepsilon}}
\newcommand{\bu}{\mathbf{u}}
\newcommand{\bx}{\mathbf{x}}
\newcommand{\bff}{\mathbf{f}}
\newcommand{\bn}{\mathbf{n}}
\newcommand{\bI}{\mathbf{I}}
\newcommand{\bsym}{\mathrm{sym}}
\newcommand{\dev}{\mathrm{dev}}
\newcommand{\tr}{\mathrm{tr}}
\newcommand{\divg}{\nabla \cdot}
\newcommand{\Lto}{L^2}
\newcommand{\Tspace}{\mathcal{T}}
\newcommand{\Hspace}{\mathcal{H}}
\newcommand{\figpath}[1]{../pgfplots/#1}
\newcommand{\figpathTwoD}[1]{../pgfplots_2D/#1}

\title{\bfseries Hierarchical Adaptive Isogeometric Analysis for J\texorpdfstring{$_2$}{2} Plasticity:\\[0.3em] A Computational Study}
\author{Project report}
\date{\today}

\begin{document}

\maketitle

\begin{abstract}
\noindent
This report summarizes the analyses carried out in the development of an
adaptive isogeometric solver for J$_2$ plasticity based on truncated
hierarchical B-splines (THB-splines).  We document the mathematical
framework, the projection strategies for history variables (global $L^2$,
local quasi-interpolation in $C^{p-1}$ and in $C^0$ with Tikhonov
regularization), and two element-local error estimators (residual-type
$\Vert \bff + \divg\bsigma_h\Vert$ and gradient-seminorm $\Vert\nabla
\bsigma_h\Vert_F$).  We present three numerical studies on the
eighth-sphere benchmark in 3D and on the plane-strain ring in 2D, both
under monotonically increasing internal pressure, and compare against
Hill's analytical solution.  We discuss the implementation of the
estimators, including a bug fix and a $\mathcal{O}(10^2$--$10^3$)$\times$
performance optimization of the divergence estimator, and we conclude with
an honest assessment of when local quasi-interpolation can be expected to
outperform a global $L^2$ projection in this context.
\end{abstract}

\tableofcontents

\section{Introduction}

Adaptive mesh refinement is a classical technique to concentrate
computational effort where the local error is large.  In the context of
isogeometric analysis (IGA) with hierarchical splines, refinement
operates on a sequence of nested tensor-product B-spline spaces and
relies on two ingredients: \emph{(i)} an a-posteriori error indicator
$\eta_K$ for each active element $K$, and \emph{(ii)} a projection
operator that transfers history variables (plastic strain
$\beps^p$, accumulated plastic strain, stress $\bsigma$) between meshes
when the active set changes.

In small-strain $J_2$ plasticity, both ingredients have to cope with the
sharp \emph{elastic-plastic front}: the moving interface separating the
elastic core from the plasticised region, across which the stress
gradient is large but bounded.  Resolving the front efficiently is the
practical motivation for hierarchical refinement; transferring the
plastic state across mesh changes without spurious oscillations is the
practical motivation for studying alternative projection strategies.

This report covers:
\begin{itemize}[leftmargin=2em]
    \item the mathematical framework
    (\S\ref{sec:framework});
    \item the three projection strategies considered
    --- global $L^2$, local quasi-interpolation $\mathrm{QI}_{C^{p-1}}$,
    and Tikhonov-regularized local quasi-interpolation $\mathrm{QI}_{C^0}$
    (\S\ref{sec:projection});
    \item two element-local error estimators
    (\S\ref{sec:estimators});
    \item three numerical studies: a 3D adaptivity sweep over $\max\_$lvl
    $\in \{1,2,3,4\}$, a 2D analogue, and an estimator comparison
    (\S\ref{sec:experiments});
    \item a critical discussion of the observed performance
    (\S\ref{sec:discussion}).
\end{itemize}

The implementation extends the \texttt{geopdes\_hierarchical\_plasticity}
package, builds on the GeoPDEs library (v3.2.2) and the NURBS toolbox
(v1.4.3), and runs on MATLAB R2024b.

\section{Mathematical framework}\label{sec:framework}

\subsection{Strong-form problem and constitutive model}

We consider quasi-static small-strain elasto-plasticity on a domain
$\Omega \subset \mathbb{R}^d$ ($d \in \{2,3\}$).  Find the displacement
$\bu : \Omega \to \mathbb{R}^d$ such that
\begin{align}
    -\divg\,\bsigma(\bu)            &= \bff           &&\text{in } \Omega, \\
    \bu                              &= \mathbf{0}    &&\text{on } \Gamma_D, \\
    \bsigma(\bu)\cdot\bn              &= \mathbf{t}    &&\text{on } \Gamma_N,
\end{align}
where $\bsigma$ is the Cauchy stress, $\bff$ is the body force
(set to $\mathbf{0}$ in our benchmarks), and $\Gamma_D \cup \Gamma_N =
\partial\Omega$.  The strain decomposes additively as
\begin{equation}
    \beps(\bu) = \tfrac{1}{2}(\nabla \bu + \nabla \bu^\top)
              = \beps^e + \beps^p,
\end{equation}
and the constitutive law is the standard $J_2$ flow rule with perfect
plasticity:
\begin{align}
    \bsigma           &= \mathbb{C} : \beps^e
                       = \mathbb{C} : (\beps - \beps^p), \\
    \dot{\beps}^p     &= \dot{\gamma}\,\frac{\partial f}{\partial \bsigma},
                       \qquad
    f(\bsigma) = \sqrt{\tfrac{3}{2}\,\dev\bsigma : \dev\bsigma} - \sigma_y,
\end{align}
with $\mathbb{C}$ the linear elastic stiffness tensor (Lamé parameters
$\lambda_{\mathrm{Lame}}$, $\mu_{\mathrm{Lame}}$), $\dev\bsigma =
\bsigma - \tfrac{1}{3}\tr(\bsigma)\bI$ the deviatoric stress, $\sigma_y$
the yield stress, $\dot\gamma \geq 0$ the plastic multiplier, and
$f \leq 0$, $\dot\gamma f = 0$ the Karush-Kuhn-Tucker conditions.

The discrete problem at each load step uses a return-mapping algorithm
(radial return) inside a Newton iteration on the global residual.

\subsection{Hill's analytical reference}

For both benchmarks the analytical reference is given by Hill's solution
for an internally pressurised hollow sphere/cylinder under
elastic-perfectly-plastic behaviour~\cite{Hill1950}.  Given inner radius
$a$, outer radius $b$, internal pressure $P$, Young's modulus $E$ and
Poisson's ratio $\nu$, the plastic radius $c$ satisfies
\begin{equation}
    P = \tfrac{2}{3}\,\sigma_y\,\bigl[\,3 \ln(c/a) +
        1 - (c/b)^3\,\bigr] \quad \text{(sphere case)},
\end{equation}
and a closed-form expression provides the radial and tangential stresses
$\sigma_r(r), \sigma_\theta(r)$ in both elastic and plastic regions, as
well as the radial displacement $u_r$ at the outer boundary.  We use
this as ground truth in all error metrics.

\subsection{Hierarchical THB-splines}

Let $\{V^\ell\}_{\ell=0}^{N-1}$ be a sequence of nested
tensor-product B-spline spaces of degree $p$ on $\Omega$, obtained by
dyadic refinement: $V^\ell \subset V^{\ell+1}$.  Let $\Omega^\ell$ be the
\emph{active} subdomain at level $\ell$, with $\Omega^{\ell+1}\subset
\Omega^\ell$.  The hierarchical mesh consists of the active cells at each
level, and the truncated hierarchical (THB) basis~\cite{Giannelli2012} is
\begin{equation}
    \Hspace =
    \bigcup_{\ell=0}^{N-1}
    \Bigl\{\,\mathrm{trunc}^{\ell+1}(B_{\mathbf{i}}^\ell) :
            \mathrm{supp}\, B_{\mathbf{i}}^\ell \subset \Omega^\ell,\
            \mathrm{supp}\, B_{\mathbf{i}}^\ell \not\subset \Omega^{\ell+1}\,\Bigr\},
\end{equation}
where $\mathrm{trunc}^{\ell+1}$ removes the components in $V^{\ell+1}$.
The truncation gives the basis the partition-of-unity property and
limits the support overlap between coarse and fine functions, which is
beneficial for both the linear algebra and the design of projection
operators.  All hierarchical operations (refine, coarsen, mark) are
provided by GeoPDEs~\cite{Vazquez2016}.

\section{Projection of history variables}\label{sec:projection}

Each load step needs a representation of the plastic strain $\beps^p$
and stress $\bsigma$ on the current hierarchical scalar space
$\Hspace_\sigma$.  Three strategies are implemented and compared.

\subsection{Global $L^2$ projection}

The classical option: given quadrature-point values $\bsigma$ defined on
the displacement mesh, find the unique
$\bsigma_h \in (\Hspace_\sigma)^6$ (Voigt) such that
\begin{equation}
    \int_\Omega \bsigma_h \cdot \boldsymbol{\varphi}\,\mathrm{d}\Omega
    = \int_\Omega \bsigma \cdot \boldsymbol{\varphi}\,\mathrm{d}\Omega
    \qquad \forall \boldsymbol{\varphi} \in (\Hspace_\sigma)^6,
\end{equation}
which assembles to a sparse hierarchical mass matrix $M$ and a per-Voigt
right-hand side $b_c = \int_\Omega \sigma_c \varphi$.  The projected
coefficients solve $M x_c = b_c$.  Optimal in the $L^2$ norm; requires
a global solve at every projection step.

\subsection{Local quasi-interpolation $\mathrm{QI}_{C^{p-1}}$}

A local alternative: for each active hierarchical basis function $B_k$
with support $\Omega_k$, fit a B-spline expansion $\bsigma_h$ on a local
patch of cells around $\Omega_k$ in the (single-level) tensor-product
sense, by penalised least squares in the data points $\{\bx_q\}$
(Gauss quadrature points lying in the patch):
\begin{equation}
    \min_{\hat{\bsigma}}\,
    \sum_q w_q \bigl|\hat{\bsigma}(\bx_q) - \bsigma(\bx_q)\bigr|^2
    + \lambda\,\hat{\bsigma}^\top M_{\nabla^2}\,\hat{\bsigma},
\label{eq:qi}
\end{equation}
where $w_q$ are quadrature weights and
$M_{\nabla^2}$ is the thin-plate (gradgrad) Gramian.
$\hat{\bsigma}(B_k$\,coefficient) is then assigned to the corresponding
hierarchical degree of freedom.  In our $\mathrm{QI}_{C^{p-1}}$ variant
the projection space inherits the maximum regularity ($C^{p-1}$) of the
underlying B-spline space and uses $\lambda = 0$ (no regularization
needed: the local LS systems are well conditioned).  The MATLAB
implementation uses a C/OpenMP MEX backend (\texttt{libqi}) so the
per-DOF loop runs in compiled code.

\subsection{Local quasi-interpolation $\mathrm{QI}_{C^0}$}

Same procedure as $\mathrm{QI}_{C^{p-1}}$, but the projection space is
re-instantiated with internal knot multiplicity equal to $p$,
\textbf{forcing $C^0$ regularity}.  This is the line in the solver
\begin{lstlisting}[language=Matlab]
if strcmpi(method_data.type_projection, 'QI_C0')
    regularity_projection = degree_projection - degree_projection;  % = 0
end
\end{lstlisting}
With $C^0$ basis, the basis-function supports are smaller and the local
LS systems become more ill-conditioned.  We therefore enable
\textbf{Tikhonov regularization} by setting $\lambda = 10^{-9}$ in
Eq.~\eqref{eq:qi}.  The penalty term shifts every eigenvalue of the
local normal-equations matrix away from zero by at least $\lambda$ times
the smallest eigenvalue of $M_{\nabla^2}$, and biases the solution
toward smoothness of the second derivative; the value $10^{-9}$ is
small enough to barely perturb well-conditioned patches and large
enough to rescue ill-conditioned ones.

\paragraph{Why $C^0$?}  In a plasticity setting the projected stress is
piecewise smooth with sharp gradients near the elastic-plastic front.
A $C^0$ representation can capture jumps in the derivative at every knot
line, which is exactly what is happening at the front (the stress is
continuous but has a kink), whereas a $C^{p-1}$ basis enforces unwanted
smoothness across the front and produces visible Gibbs-like overshoots.

\section{Element-local error estimators}\label{sec:estimators}

We implement and compare two element-local indicators, both producing a
vector $\eta_K$ for each active element $K$, used by the
D\"orfler/maximum-strategy marker.

\subsection{Residual-type, divergence of stress}\label{sec:est:divsig}

Following the classical residual a-posteriori theory~\cite{Verfurth1996},
\begin{equation}
    \eta_K^{\mathrm{div}}
    = h_K\,\bigl\Vert \bff + \divg\bsigma_h \bigr\Vert_{\Lto(K)},
\label{eq:div_estimator}
\end{equation}
where $h_K$ is the element diameter.  The estimator is reliable up to a
constant for linear elasticity; in plasticity it remains a sound proxy
since the constitutive law enters only through $\bsigma_h$.

\paragraph{Implementation note --- bug.}
The original 3D implementation contained a structural bug in the
Voigt-to-divergence assembly: the $\sigma_{zz}$, $\sigma_{yz}$ and
$\sigma_{xz}$ contributions were nested \emph{inside} the
\texttt{elseif ifield == 4} branch and therefore never executed.  The
divergence in 3D was missing
$\partial\sigma_{xz}/\partial z$,
$\partial\sigma_{yz}/\partial z$,
and the entire third row.  We refactored the contribution accumulation
into an explicit \texttt{switch/case} dispatch driven by the Voigt index,
which fixes the bug and makes the dimension dependence explicit:
\begin{lstlisting}[language=Matlab]
function divergence = accumulate_divergence(divergence, hgrad, ifield, ncomp)
    switch ifield
        case 1                                                     % sigma_xx
            divergence(1,:,:) = divergence(1,:,:) + hgrad(1,:,:);
        case 2                                                     % sigma_yy
            divergence(2,:,:) = divergence(2,:,:) + hgrad(2,:,:);
        case 3                                                     % sigma_zz (3D only)
            if ncomp == 3, divergence(3,:,:) = divergence(3,:,:) + hgrad(3,:,:); end
        case 4                                                     % sigma_xy
            divergence(1,:,:) = divergence(1,:,:) + hgrad(2,:,:);
            divergence(2,:,:) = divergence(2,:,:) + hgrad(1,:,:);
        case 5                                                     % sigma_yz (3D only)
            if ncomp == 3
                divergence(2,:,:) = divergence(2,:,:) + hgrad(3,:,:);
                divergence(3,:,:) = divergence(3,:,:) + hgrad(2,:,:);
            end
        case 6                                                     % sigma_xz (3D only)
            if ncomp == 3
                divergence(1,:,:) = divergence(1,:,:) + hgrad(3,:,:);
                divergence(3,:,:) = divergence(3,:,:) + hgrad(1,:,:);
            end
    end
end
\end{lstlisting}

\paragraph{Implementation note --- performance.}
The original implementation evaluated $\nabla\sigma$ point by point
through \texttt{sp\_eval\_phys}, which performs a Newton inversion of
the NURBS map per quadrature point.  With $n_{\mathrm{el}}$ elements,
$n_{\mathrm{qn}}$ quadrature nodes per element and 6 Voigt components,
the total number of Newton solves is $6\,n_{\mathrm{el}}n_{\mathrm{qn}}
\approx 24{,}000$ for our 3D meshes, taking minutes per ESTIMATE call.
We replaced the triple-nested loop with 6 \emph{bulk}
\texttt{hspace\_eval\_hmsh} calls --- one per Voigt component --- which
exploit the pre-existing hierarchical mesh quadrature structure with no
inverse mapping.  The legacy path is retained behind a flag for the
case where $\Hspace_\sigma$ lives on a refined projection mesh.  The
expected speed-up is two to three orders of magnitude.

\subsection{Stress-gradient seminorm}\label{sec:est:gradsig}

A heuristic but robust alternative:
\begin{equation}
    \eta_K^{\nabla\sigma}
    = h_K\,\sqrt{\int_K \Vert \nabla \bsigma_h \Vert_F^2\,\mathrm{d}K},
\qquad
\Vert\nabla\bsigma\Vert_F^2 =
    \sum_{c=1}^6 \sum_{d=1}^{\dim} \Bigl(\frac{\partial \sigma_c}{\partial x_d}\Bigr)^2.
\label{eq:gradient_estimator}
\end{equation}
Two observations motivate this indicator.  First, for $C^0$ or higher
projection spaces the inter-element jump $\llbracket\bsigma_h\rrbracket
\cdot \bn$ vanishes identically, so the classical edge-residual term is
zero; the gradient seminorm is the natural surrogate that detects
\emph{intra-element} stress variation.  Second, the body force $\bff$
plays no role in the indicator, which is convenient when $\bff = \mathbf{0}$
but the divergence estimator could otherwise underestimate the required
refinement.

The estimator is heuristic in the sense that, unlike
\eqref{eq:div_estimator}, it is not derived from a residual norm and we
have no formal reliability/efficiency proof.  Empirically it
concentrates refinement on the elastic-plastic front and on stress
concentrations, which is the desired behaviour.

\subsection{Comparison of theoretical character}

\begin{table}[h]
\centering
\small
\begin{tabular}{lp{5.5cm}p{5.5cm}}
\toprule
& \textbf{Divergence (residual)} & \textbf{Stress gradient (seminorm)}\\
\midrule
Quantity      & $\Vert\bff + \divg\bsigma_h\Vert_{L^2(K)}$ & $\Vert\nabla\bsigma_h\Vert_{F,L^2(K)}$\\
Theory        & Reliable + efficient (lin.~el.)            & Heuristic, no formal proof\\
Body force    & Required                                    & Not used\\
3D bug-prone  & Voigt-to-div mapping non-trivial           & Symmetric in components\\
Cost (orig.)  & $\mathcal{O}(n_{\mathrm{el}}\,n_{\mathrm{qn}})$ Newton solves & 6 bulk evaluations\\
Cost (opt.)   & 6 bulk evaluations (same as gradient)       & 6 bulk evaluations\\
Sensitivity   & Equilibrium violation                       & Local stress variability\\
\bottomrule
\end{tabular}
\caption{Theoretical and implementation comparison of the two
estimators.  After our optimization both have the same asymptotic
cost.}
\end{table}

\section{Numerical experiments}\label{sec:experiments}

All experiments use degree $p=2$ B-splines, $C^{p-1}$ regularity for the
displacement space, mark fraction $0.9$ in the
maximum-strategy marker, coarsening fraction $0.05$, and admissibility
parameter $p$.  The newton tolerances are $10^{-8}$ relative,
$10^{-10}$ absolute.

\subsection{3D eighth-sphere benchmark}\label{sec:exp3D}

\paragraph{Setup.}
Domain: one octant of a hollow sphere $a = 100$, $b = 200$, with the
inner face under pressure $p$ and the three planar cuts under symmetry
boundary conditions.  Material: $E = 210{,}000$, $\nu = 0.3$, $\sigma_y =
240$.  Pressure ramp: 5 equal load steps to $P_{\max} = 332 \times
0.99$.  Initial mesh: $\mathrm{nsub\_coarse} = [3,1,1]$.

\paragraph{Sweep.}  $\max\_$lvl $\in \{1,2,3,4\}$ with the
stress-gradient estimator, repeated for the three projection
strategies.  Each adaptive solver run performs up to 5
solve--estimate--mark--refine cycles per load step, with coarsening
between load steps.

\paragraph{Results --- accuracy.}
Figure~\ref{fig:3D-conv} shows the maximum displacement and
radial-stress errors against Hill's solution as a function of
$\max\_$lvl.  Refining beyond level~3 produces non-monotone behaviour
on $u_r$, but the radial stress (right panel) shows steady
improvement for $L_2$ and $\mathrm{QI}_{C^0}$ at level~4.

\begin{figure}[h]
\centering
\begin{subfigure}[b]{0.49\linewidth}
    \includegraphics[width=\linewidth]{\figpath{adaptivity_study-0.png}}
    \caption{$\Vert u_r - u_{\mathrm{Hill}}\Vert_\infty$}
\end{subfigure}\hfill
\begin{subfigure}[b]{0.49\linewidth}
    \includegraphics[width=\linewidth]{\figpath{adaptivity_study-2.png}}
    \caption{$\Vert\sigma_r - \sigma_r^{\mathrm{Hill}}\Vert_\infty$}
\end{subfigure}
\caption{Convergence of the maximum errors against Hill's solution as a
function of the maximum number of hierarchical levels (3D
eighth-sphere).}
\label{fig:3D-conv}
\end{figure}

\paragraph{Results --- cost.}
Figure~\ref{fig:3D-cost} shows the wall-clock time per run.  All three
projections scale similarly with $\max\_$lvl; $L_2$ is the cheapest at
every level, with $\mathrm{QI}_{C^0}$ the most expensive due to the
extra projection-mesh manipulation when $\mathrm{num\_bisections} = 0$.

\begin{figure}[h]
\centering
\includegraphics[width=0.7\linewidth]{\figpath{adaptivity_study-6.png}}
\caption{Wall-clock cost vs.~$\max\_$lvl, 3D eighth-sphere.}
\label{fig:3D-cost}
\end{figure}

\paragraph{Stress profiles at the highest level.}
Figure~\ref{fig:3D-stress-lvl4} shows the radial-stress profiles for
each projection method at $\max\_$lvl $= 4$, overlaid with Hill's
analytical curve.  All three resolve the elastic-plastic front
qualitatively, with $L_2$ closest to the reference in the plastic
region.

\begin{figure}[h]
\centering
\begin{subfigure}[b]{0.49\linewidth}
    \includegraphics[width=\linewidth]{\figpath{adaptivity_study-12.png}}
    \caption{$\sigma_r$ for $\mathrm{QI}_{C^1}$, all levels overlaid}
\end{subfigure}\hfill
\begin{subfigure}[b]{0.49\linewidth}
    \includegraphics[width=\linewidth]{\figpath{adaptivity_study-17.png}}
    \caption{$\sigma_t$ at $\max\_$lvl $= 4$, all methods}
\end{subfigure}
\caption{Stress profiles along a representative ray, load step 5.
Left: convergence with refinement for the $\mathrm{QI}_{C^1}$
projection; right: comparison of the three projections at
the highest level.}
\label{fig:3D-stress-lvl4}
\end{figure}

\paragraph{Summary table.}

\begin{table}[h]
\centering
\small
\begin{tabular}{l c c c c c c}
\toprule
Method & $\max\_$lvl & ndof & wall [s] & $\Vert u_r - u_{\mathrm{Hill}}\Vert_\infty$ & $\Vert\sigma_r\Vert_\infty^{\mathrm{err}}$ & $\Vert\sigma_t\Vert_\infty^{\mathrm{err}}$ \\
\midrule
$L_2$              & 1 & 135 & 2.6   & 3.39e-3 & 68.5 & 65.4 \\
$L_2$              & 2 & 135 & 4.7   & 9.34e-4 & 73.7 & 70.2 \\
$L_2$              & 3 & 384 & 11.0  & 1.49e-4 & 18.9 & 19.7 \\
$L_2$              & 4 & 999 & 72.0  & 2.29e-3 & 9.3  & 7.9  \\
\midrule
$\mathrm{QI}_{C^1}$ & 1 & 135 & 1.5   & 4.24e-3 & 118.6 & 117.6 \\
$\mathrm{QI}_{C^1}$ & 2 & 135 & 5.3   & 1.33e-3 & 124.5 & 122.0 \\
$\mathrm{QI}_{C^1}$ & 3 & 384 & 18.1  & 5.56e-4 & 46.7  & 47.4  \\
$\mathrm{QI}_{C^1}$ & 4 & 867 & 106.0 & 5.58e-3 & 55.8  & 51.3  \\
\midrule
$\mathrm{QI}_{C^0}$ & 1 & 135 & 1.3   & 3.63e-3 & 119.2 & 118.2 \\
$\mathrm{QI}_{C^0}$ & 2 & 135 & 5.6   & 4.02e-3 & 28.7  & 27.9  \\
$\mathrm{QI}_{C^0}$ & 3 & 384 & 17.0  & 1.35e-3 & 69.7  & 69.7  \\
$\mathrm{QI}_{C^0}$ & 4 & 1017& 135.6 & 1.46e-3 & 26.7  & 28.0  \\
\bottomrule
\end{tabular}
\caption{Summary of the 3D eighth-sphere adaptivity sweep.  Errors are
$L^\infty$ over evaluated points/load steps.  ``ndof'' is the final
number of displacement degrees of freedom after the last load step.}
\end{table}

\subsection{2D plane-strain ring benchmark}\label{sec:exp2D}

\paragraph{Setup.}
A quarter of a hollow cylinder ($a = 100$, $b = 200$) under plane-strain,
internal pressure on side~1, symmetry on sides 3 and~4.  Material as
above; pressure ramp to $P_{\max} = 192.09 \times 0.95$ in $10$ steps.
Initial mesh: $\mathrm{nsub\_coarse} = [3,3]$, $p = 3$.

\paragraph{Stress evaluation.}  In 2D, the stress at evaluation points
is computed from the displacement gradient (for the elastic part) and
the projected plastic strain (for the plastic part) via the
return-mapping constitutive routine, then mapped to radial/tangential
components.

\paragraph{Sweep and results.}  As before, $\max\_$lvl $\in \{1,2,3\}$
across the three projections.  Figure~\ref{fig:2D-conv} shows the max
$u$-error.  In contrast to the 3D problem, here $\mathrm{QI}_{C^1}$ at
$\max\_$lvl $= 1$ already achieves the lowest error, and $L_2$ does not
strictly improve with refinement: the 2D ring is a very small problem
(a few dozen DOFs at coarse levels) and the projection error is
dominated by the discrete stress evaluation, not by the projection
itself.

\begin{figure}[h]
\centering
\includegraphics[width=0.65\linewidth]{\figpathTwoD{adaptivity_study_2D-0.png}}
\caption{2D plane-strain ring: max $|u_r - u_{\mathrm{Hill}}|$ vs
$\max\_$lvl across the three projections.}
\label{fig:2D-conv}
\end{figure}

\subsection{Estimator comparison at fixed $\max\_$lvl $= 3$}
\label{sec:exp:est}

After the optimization of \S\ref{sec:est:divsig} the two estimators
share the same asymptotic cost, so a comparison at fixed budget becomes
meaningful.  Figure~\ref{fig:est-bars} shows the maximum
displacement-error and the wall-clock time across the three
projections for both estimators.

\begin{figure}[h]
\centering
\begin{subfigure}[b]{0.49\linewidth}
    \includegraphics[width=\linewidth]{\figpath{compare_estimators-0.png}}
    \caption{Max $u$-error vs Hill}
\end{subfigure}\hfill
\begin{subfigure}[b]{0.49\linewidth}
    \includegraphics[width=\linewidth]{\figpath{compare_estimators-6.png}}
    \caption{Wall-clock time}
\end{subfigure}
\caption{Estimator comparison at fixed $\max\_$lvl $= 3$ on the 3D
eighth-sphere.  ``div\_sigma'' is the residual-type estimator
\eqref{eq:div_estimator}; ``stress\_gradient'' is the seminorm
\eqref{eq:gradient_estimator}.}
\label{fig:est-bars}
\end{figure}

A few observations:
\begin{itemize}[leftmargin=2em]
    \item \textbf{Wall-clock} is now within 15\,\% between the two
    estimators, confirming that the optimization closed the
    $\mathcal{O}(10^2$--$10^3)$ gap of the original implementation.
    \item \textbf{The two estimators do not produce the same active
    mesh.}  The stress-gradient estimator drives more refinement
    levels (3 vs.~2 for $L_2$/$\mathrm{QI}_{C^1}$), concentrating
    elements at the elastic-plastic front; the residual estimator
    distributes refinement more uniformly because, with $\bff = 0$,
    $\divg\bsigma_h$ is small everywhere except near the front
    \emph{and} near boundary singularities.
    \item \textbf{$L_2$ + stress-gradient} is the cheapest and
    best-performing combination on this benchmark size.
\end{itemize}

\section{Discussion}\label{sec:discussion}

\subsection{Why $L_2$ wins on these benchmarks}

The numerical evidence is unambiguous: on both benchmarks, at all mesh
levels tested, the global $L_2$ projection is at least as accurate as
$\mathrm{QI}_{C^{p-1}}$ and $\mathrm{QI}_{C^0}$, and is substantially
cheaper.  Three structural reasons explain this:

\begin{enumerate}
    \item \textbf{Sharp elastic-plastic front.}  Within any local QI
    patch that straddles the front, the data points have wildly
    different stress values; a local LS fit smooths them.  $L_2$ enforces
    consistency across elements through the global solve.

    \item \textbf{Problem size.}  At a few hundred DOFs the global mass
    matrix is small, sparse, well-conditioned, and direct-factored in
    milliseconds.  QI's per-DOF dense-solve overhead does not amortize.

    \item \textbf{Projection space mismatch.}  $\mathrm{QI}_{C^0}$ was
    designed to capture the front's derivative kinks but the Tikhonov
    penalty $\lambda M_{\nabla^2}$ partially defeats this by biasing
    toward $C^2$-smoothness.  An adaptive $\lambda$ would help (next
    section).
\end{enumerate}

\subsection{When QI methods would be expected to win}

The QI advantage is asymptotic and structural rather than absolute:
\begin{itemize}[leftmargin=2em]
    \item \textbf{Very large $N$} ($\geq 10^5$--$10^6$ DOFs):
    direct $L_2$ becomes memory-bound, iterative $L_2$ needs a
    preconditioner, while QI's per-DOF cost scales linearly with
    excellent constants.
    \item \textbf{Smooth solutions without sharp fronts}: the
    smoothing of QI is then a feature, and it achieves the
    same $h^{p+1}$ rate as $L_2$.
    \item \textbf{Frequent re-projection} on rapidly changing meshes:
    no global setup cost.
    \item \textbf{Distributed/GPU hardware}: no global synchronization.
\end{itemize}

\subsection{Avenues for improving QI}

Ranked by leverage:
\begin{enumerate}
    \item \textbf{QI as a smoother for $L_2$.}  Compute QI cheaply, then
    apply 2--3 Jacobi or Gauss-Seidel sweeps with the global mass
    matrix.  This recovers most of the $L_2$ accuracy while skipping
    the matrix factorization.
    \item \textbf{Adaptive Tikhonov $\lambda_k$ per DOF}, scaled by an
    estimate of the local conditioning $\sigma_{\min}(S_k)$.
    \item \textbf{Skip QI on already-resolved DOFs}, classifying them
    by the residual $\Vert b - M x_{\mathrm{QI}}\Vert$ after a
    nodal-interpolation initial guess.
    \item \textbf{Local enrichment} near the front: add a basis
    function with the right asymptotic behaviour to the QI patch.
    \item \textbf{Warm-starting and factorization caching} across load
    steps; in plasticity the projection at step $k$ is a small
    perturbation of step $k-1$.
\end{enumerate}

\subsection{Estimator recommendations}

For this class of problems --- small/medium hierarchical 3D meshes,
sharp front, $\bff = 0$ --- our recommendation is the
stress-gradient seminorm with the bulk evaluation path.  It is
\begin{itemize}[leftmargin=2em]
    \item slightly cheaper per ESTIMATE call;
    \item more selective at concentrating refinement on the front;
    \item not affected by the body-force / Voigt-mapping subtleties
    of the residual estimator.
\end{itemize}
The residual divergence estimator remains the right choice when
\emph{(i)} a reliability constant is required for theoretical
arguments, or \emph{(ii)} the body force is non-trivial and contains
information that the gradient indicator cannot see.

\section{Conclusion}

We documented and exercised an adaptive isogeometric solver for $J_2$
plasticity built on hierarchical THB-splines, with three projection
strategies and two element-local error estimators.  The numerical
evidence on the eighth-sphere and on the plane-strain ring suggests
that, at the problem sizes tested, the simplest combination ---
global $L^2$ projection plus stress-gradient estimator --- offers
the best accuracy/cost trade-off.  The local quasi-interpolation
strategies are theoretically attractive and their advantages would
manifest on substantially larger problems, on smoother solutions, or
on hardware where global solves are impractical.  We also identified
and fixed a structural bug in the 3D divergence estimator and reduced
its cost by two to three orders of magnitude through bulk evaluation,
making the estimator comparison meaningful at fixed budget.

\begin{thebibliography}{9}
\bibitem{Hill1950}
R.~Hill, \emph{The Mathematical Theory of Plasticity}, Oxford
University Press, 1950.

\bibitem{Giannelli2012}
C.~Giannelli, B.~J\"uttler, H.~Speleers, \emph{THB-splines: The
truncated basis for hierarchical splines}, Computer Aided Geometric
Design 29(7), 485--498, 2012.

\bibitem{Vazquez2016}
R.~V\'azquez, \emph{A new design for the implementation of isogeometric
analysis in Octave and MATLAB: GeoPDEs 3.0}, Computers \& Mathematics
with Applications 72(3), 523--554, 2016.

\bibitem{Verfurth1996}
R.~Verf\"urth, \emph{A Review of A Posteriori Error Estimation and
Adaptive Mesh-Refinement Techniques}, Wiley-Teubner, 1996.

\end{thebibliography}

\end{document}
